\documentclass[journal]{IEEEtran}

% ---------------------------------------------------------------
% Packages
% ---------------------------------------------------------------
\usepackage[T1]{fontenc}
\usepackage[utf8]{inputenc}
\usepackage{amsmath,amssymb,amsfonts}
\usepackage{algorithmic}
\usepackage{graphicx}
\usepackage{textcomp}
\usepackage{xcolor}
\usepackage{booktabs}
\usepackage{multirow}
\usepackage{url}
\usepackage{cite}
\usepackage{caption}
\usepackage{subcaption}
\usepackage{array}
\usepackage{tabularx}
\usepackage{makecell}
\usepackage{hyperref}

\hypersetup{
    colorlinks=true,
    linkcolor=blue,
    citecolor=blue,
    urlcolor=blue
}

% ---------------------------------------------------------------
% Title and Authors
% ---------------------------------------------------------------
\begin{document}

\title{NutriBiteBot: A TabNet-Driven Framework for Ingredient-Level
Dietary Portion Recommendation in Patients with Hypertension,
Type~2 Diabetes Mellitus, and Chronic Kidney Disease}

\author{
    \IEEEauthorblockN{Mohana\IEEEauthorrefmark{1},
                      B.~G.~Sudarshan\IEEEauthorrefmark{1},
                      Mrida~Pradhan\IEEEauthorrefmark{1}, and
                      Gayathri~Sunil~Nambiar\IEEEauthorrefmark{1}}
    \IEEEauthorblockA{\IEEEauthorrefmark{1}Department of Computer Science and Engineering,
    RV College of Engineering, Bengaluru, Karnataka 560059, India\\
    \{mohana, sudharshanbg, mridapradhan.cd23, gayathrisuniln.cd23\}@rvce.edu.in}
}

\maketitle

% ---------------------------------------------------------------
% Abstract
% ---------------------------------------------------------------
\begin{abstract}
Patients simultaneously diagnosed with hypertension (HTN), Type~2 Diabetes
Mellitus (T2DM), and Chronic Kidney Disease (CKD) face irreconcilable
dietary constraints that no existing clinical nutrition tool resolves
automatically. We present \textbf{NutriBiteBot}, an interpretable,
end-to-end clinical decision support system that computes
ingredient-level dietary portion recommendations directly from 14
patient laboratory and clinical features. The system employs TabNet
classifiers with monotonic feature transformations for per-nutrient
risk severity inference, a five-constraint sigmoid-weighted portion
computation engine grounded in the Indian Food Composition Tables (IFCT
2017), a Roboflow computer-vision ingredient detector, a
KDIGO-compliant constraint relaxation hierarchy, and a Groq Llama-3.3
generative agent that operates exclusively within hard gram-level
boundaries set by the upstream rules engine. All components execute
within a local Docker environment to satisfy India's Digital Personal
Data Protection Act (DPDPA) 2023. Our TabNet clinical classifier
achieves a mean accuracy of 99.91\% and mean weighted F1 of 99.91\%
across four sensitivity targets (sodium, potassium, protein, and
carbohydrate), with mean Cohen's $\kappa = 0.9983$. We compare TabNet
against NGBoost as a competitive gradient-based baseline and demonstrate
that TabNet's attention-based feature attribution, native
interpretability, and compatibility with monotonic preprocessing
constraints make it the preferred architecture for clinical tabular
inference in this setting. NutriBiteBot constitutes the first system to
combine real-time computer vision, ML-derived clinical constraints,
LLM-bounded recipe generation, and on-premises data persistence for
multi-comorbid dietary management in an Indian clinical context.
\end{abstract}

\begin{IEEEkeywords}
Clinical decision support, dietary recommendation, TabNet, NGBoost,
chronic kidney disease, type 2 diabetes, hypertension, Indian food
composition, IFCT 2017, KDIGO, interpretable machine learning,
tabular deep learning, ingredient-level nutrition
\end{IEEEkeywords}

% ---------------------------------------------------------------
\section{Introduction}
\label{sec:intro}
% ---------------------------------------------------------------

India bears one of the highest burdens of non-communicable disease
(NCD) globally. As of 2024, over 101~million Indians live with Type~2
Diabetes Mellitus (T2DM), approximately 220~million with hypertension
(HTN), and an estimated 17.2\% of the adult population shows markers of
Chronic Kidney Disease (CKD)~\cite{anjana2023}. These conditions
frequently co-occur in the same patient, constituting a clinical triad
that imposes extreme and often contradictory demands on dietary
management.

Diet is a primary therapeutic lever for all three conditions, yet the
governing clinical guidelines conflict at the nutrient level in ways
that are non-trivial to resolve without expert knowledge. The DASH diet
for HTN prescribes high potassium ($\geq$4,700~mg/day); CKD management
mandates strict potassium restriction ($\leq$2,000~mg/day) to prevent
fatal hyperkalemia. Adequate dietary protein supports glycaemic control
in T2DM, while CKD requires protein restriction (0.6--0.8~g/kg/day) to
slow renal decline~\cite{kdoqi2020,ada2024}. No existing tool resolves
these conflicts automatically, and no tool translates clinical nutrient
limits into ingredient-level gram quantities anchored to the patient's
actual laboratory values.

Current practice depends on manual dietitian consultation, a
30--60~minute process per patient that is prone to error when three or
more conflicting guideline sets must be reconciled simultaneously.
India faces a severe shortage of registered dietitians, estimated at a
national ratio of 1 per 100,000 population, a scalability gap
unaddressed by any existing clinical nutrition software. Furthermore,
most Indian households rely on ingredients catalogued in the Indian Food
Composition Tables (IFCT 2017)~\cite{ifct2017}, a database entirely
absent from international dietary tools such as MyFitnessPal, Cronometer,
or Nutritics.

This paper presents \textbf{NutriBiteBot}, a locally-deployed,
interpretable clinical decision support system that closes these gaps
through eight functional layers: multi-source data ingestion, ML model
training, per-patient risk inference with clinical explanation
generation, session-persistent nutrient budget management,
five-constraint sigmoid-weighted ingredient-level portion computation,
cosine-similarity-based substitution for contra-indicated ingredients,
caloric sufficiency validation with KDIGO-compliant constraint
relaxation, and structured output generation with daily ledger tracking.

The core contributions of this work are as follows:
\begin{itemize}
    \item A multi-target TabNet classifier that infers per-nutrient
    sensitivity levels (Low/Moderate/High) for sodium, potassium,
    protein, and carbohydrates from 14 EHR-derived clinical features,
    achieving mean accuracy 99.91\% and mean weighted F1 99.91\% with
    mean Cohen's $\kappa = 0.9983$.
    \item A systematic comparison of TabNet against NGBoost as a
    competitive probabilistic gradient-boosting baseline, motivating the
    selection of TabNet on interpretability, monotonicity compatibility,
    and attention-based feature attribution grounds.
    \item A novel five-constraint sigmoid severity-to-fraction portion
    computation formulation that produces ingredient-specific maximum safe
    gram quantities from continuous severity scores against a live daily
    nutrient budget.
    \item A KDIGO 2024-compliant hierarchical constraint relaxation
    hierarchy that prevents caloric collapse without violating potassium
    and sodium safety limits.
    \item A hybrid bounded-LLM architecture wherein a Groq Llama-3.3
    generative agent produces adapted recipes exclusively within hard
    numerical boundaries enforced by the upstream rules engine.
    \item A fully local Docker deployment that ensures patient laboratory
    data never traverses a public network, satisfying DPDPA 2023
    requirements.
\end{itemize}

The remainder of this paper is organised as follows.
Section~\ref{sec:related} reviews related work.
Section~\ref{sec:architecture} describes the system architecture.
Section~\ref{sec:mlpipeline} details the ML pipeline.
Section~\ref{sec:ngboost} presents the NGBoost baseline.
Section~\ref{sec:comparison} compares TabNet and NGBoost.
Section~\ref{sec:portion} describes the portion computation engine.
Section~\ref{sec:results} presents experimental results.
Section~\ref{sec:usecase} illustrates a clinical use case.
Section~\ref{sec:discussion} discusses limitations and future work.
Section~\ref{sec:conclusion} concludes.

% ---------------------------------------------------------------
\section{Related Work}
\label{sec:related}
% ---------------------------------------------------------------

\subsection{Clinical Dietary Decision Support Systems}
Existing clinical nutrition tools fall into two categories: consumer
macronutrient trackers (MyFitnessPal, Cronometer, Healthify Me) and
professional dietitian platforms (Nutritics, Computrition). Consumer
tools perform calorie accounting with no disease-specific constraint
logic and no mechanism for resolving inter-condition dietary conflicts.
Professional platforms provide nutrient analysis but require a
registered dietitian to manually configure every constraint and resolve
conflicts --- precisely the bottleneck that NutriBiteBot automates.
Critically, no existing tool grounds ingredient-level portion
computation in patient-specific laboratory values, and none integrates
Indian food composition data (IFCT 2017). The closest deployed system
for the Indian market, Healthify Me, routes all data through cloud
infrastructure and cannot enforce hard gram-level clinical ceilings.

\subsection{Machine Learning for Clinical Tabular Data}
The challenge of learning from tabular EHR data with heterogeneous
feature types and modest sample sizes has motivated a range of
approaches. Gradient boosting methods (XGBoost~\cite{chen2016xgboost},
LightGBM, NGBoost~\cite{duan2020ngboost}) dominate tabular benchmarks
due to their inductive bias for heterogeneous feature distributions and
resistance to overfitting on small datasets. NGBoost extends gradient
boosting to probabilistic output via natural gradient estimation,
producing per-prediction uncertainty estimates that are clinically
meaningful. TabNet~\cite{arik2021tabnet} proposes an attention-based
deep architecture that processes tabular features without
featurisation, learning sparse feature selection masks at each decision
step using a sequentially attentive architecture. TabNet has been shown
to match or exceed gradient boosting performance on several large
tabular benchmarks while producing per-instance feature attribution
masks that are native to the inference computation --- a property of
direct clinical value for generating patient-specific explanations.

\subsection{Monotonicity in Clinical ML}
Physiologically motivated monotonicity constraints have been applied in
clinical predictive models to prevent counter-intuitive inferences (e.g.,
lower predicted risk at higher creatinine). Isotonic regression,
monotone neural networks, and constrained gradient boosting have all been
applied to enforce such constraints~\cite{gupta2019}. The combination of
monotonic preprocessing with a native attention-based feature selector
--- as employed in NutriBiteBot --- has not previously been reported in
the clinical dietary recommendation literature.

\subsection{LLM-Based Dietary Advice}
Recent work has explored large language models (LLMs) for dietary
guidance~\cite{xu2024llm}. However, unconstrained LLMs are unsuited for
clinical dietary recommendation: they lack deterministic safety guarantees
and can produce clinically plausible but unsafe portion quantities.
NutriBiteBot addresses this by employing a hybrid architecture in which
Groq Llama-3.3 generates recipes only after a rules engine has
computed hard gram-level constraint boundaries that the agent cannot
exceed.

% ---------------------------------------------------------------
\section{System Architecture}
\label{sec:architecture}
% ---------------------------------------------------------------

NutriBiteBot is deployed as a locally containerised system comprising
five distinct architectural components: a single-page web frontend, a
Python/Flask backend server, a four-stage ML pipeline, a local Supabase
(PostgreSQL) database, and a Docker deployment environment
(Fig.~\ref{fig:architecture}). All components execute within a local
Docker environment; patient laboratory data, generated recipes, and
clinical records never leave the institution's infrastructure. This
design satisfies the sensitivity requirements of India's DPDPA 2023.

\begin{figure}[!t]
    \centering
    % Replace with actual figure: \includegraphics[width=\columnwidth]{figures/architecture.png}
    \fbox{\parbox{0.9\columnwidth}{\centering\vspace{2cm}
        [Figure: Complete Architecture of NutriBiteBot ---\\
        six-step data flow diagram]
    \vspace{2cm}}}
    \caption{Complete architecture of NutriBiteBot. The six-step
    data flow: (1)~patient uploads lab values and fridge image via the
    frontend; (2)~backend POSTs image to the Vision Detection module;
    (3)~Rules Engine returns safe portion numerical constraints to the
    Generative Agent; (4)~Generative Agent returns adapted recipe to
    backend; (5)~backend persists records to local Supabase; (6)~backend
    returns triage summary and recipe to frontend.}
    \label{fig:architecture}
\end{figure}

\subsection{Frontend}
The user interface is a single-page web application (HTML5/CSS3/
JavaScript) exposing three functional panels: a Fridge Scanner
(refrigerator image upload and ingredient display), a Patient Data Input
form (14 clinical features), and a Recipe Generator panel (safe portion
recommendations and adapted recipe). The frontend communicates exclusively
via REST API calls to the backend and holds no patient data in
client-side storage.

\subsection{Backend Server}
The backend is a Python/Flask server that orchestrates ML inference
across the four-stage pipeline, aggregates rules engine constraint
outputs, and manages write operations to the local database. Backend
ML inference uses PyTorch as the TabNet runtime. Agent orchestration
logic sequences pipeline stages and manages error handling, fallback
logic, and caloric sufficiency validation.

\subsection{ML/DL Pipeline}
The four-stage pipeline (Section~\ref{sec:mlpipeline}) executes
sequentially: Vision Detection (Stage~1), TabNet Risk Stratification
(Stage~2), Clinical Rules Engine with Portion Computation (Stage~3),
and Bounded Generative Agent (Stage~4).

\subsection{Database}
Patient records and generated recipes are persisted to a Supabase
instance running on local PostgreSQL within the same Docker environment.
No data is written to Supabase's cloud infrastructure. Stored entities
include patient identifiers, clinical feature vectors, per-session
portion decision records, generated recipe outputs, and daily nutrient
ledger state.

\subsection{Deployment}
The complete system is containerised via a single \texttt{docker-compose}
configuration that defines all inter-service networking. No component is
exposed to the public internet by default. This model supports hospital
on-premises installation, clinic workstations, and telemedicine platforms
operating within a secured local network.

% ---------------------------------------------------------------
\section{ML Pipeline: Model Training and Inference}
\label{sec:mlpipeline}
% ---------------------------------------------------------------

\subsection{Dataset and Label Generation}

Training data are drawn from MIMIC-IV (version~2.2)~\cite{mimic2023},
a publicly available EHR dataset. A threshold-based \texttt{LabelGenerator}
produces four sensitivity targets without requiring learned labels:

\begin{itemize}
    \item \textbf{sodium\_sensitivity}: derived from serum Na,
    HTN diagnosis flag, and systolic blood pressure (SBP).
    \item \textbf{potassium\_sensitivity}: derived from serum K and
    estimated glomerular filtration rate (eGFR).
    \item \textbf{protein\_restriction}: derived from eGFR, serum
    creatinine, and CKD diagnosis flag.
    \item \textbf{carb\_sensitivity}: derived from HbA1c, fasting blood
    sugar (FBS), BMI, and DM diagnosis flag.
\end{itemize}

Each target is classified as Low (0), Moderate (1), or High (2).
Phosphorus sensitivity is derived by rule from eGFR thresholds per
KDIGO 2017~\cite{kdigo2012}.

\subsection{Preprocessing and Monotonic Constraints}

Preprocessing applies median imputation for missing values and standard
(z-score) scaling. A \texttt{MonotonicFeatureTransformer} enforces
physiologically motivated monotonicity via isotonic regression:
creatinine, serum K, HbA1c, FBS, SBP, DBP, BMI, and serum Na are
constrained to positive monotonicity (higher value $\Rightarrow$ higher
predicted risk class); eGFR is constrained to negative monotonicity
(lower eGFR $\Rightarrow$ higher renal risk class). This prevents the
model from learning counter-physiological associations irrespective of
distributional artefacts in the training data.

\subsection{TabNet Classifier Configuration}

Four \texttt{TabNetClassifier} instances (one per target) are trained
independently. Key hyperparameters are: sparsemax attention activation,
\texttt{n\_steps}~$= 5$, \texttt{n\_a}~$= 64$, \texttt{n\_d}~$= 64$,
\texttt{momentum}~$= 0.02$, \texttt{gamma}~$= 1.3$, and
\texttt{epsilon}~$= 1 \times 10^{-15}$.
Each model outputs a \texttt{predict\_proba} vector and a per-patient
feature attribution mask. The attribution masks are used directly in
Layer~2 to generate clinical explanations that cite the patient's
specific laboratory values and the driving feature for each sensitivity
classification.

\begin{figure}[!t]
    \centering
    % Replace with: \includegraphics[width=\columnwidth]{figures/ml_pipeline.png}
    \fbox{\parbox{0.9\columnwidth}{\centering\vspace{2.5cm}
        [Figure: Four-Stage ML/DL Pipeline Diagram]
    \vspace{2.5cm}}}
    \caption{The four-stage ML/DL pipeline. Stage~1: Roboflow CV
    ingredient detection. Stage~2: TabNet multi-target risk stratification.
    Stage~3: Clinical Rules Engine with sigmoid portion computation.
    Stage~4: Groq Llama-3.3 bounded generative agent.}
    \label{fig:mlpipeline}
\end{figure}

% ---------------------------------------------------------------
\section{NGBoost Baseline}
\label{sec:ngboost}
% ---------------------------------------------------------------

To establish a competitive probabilistic baseline and motivate the
selection of TabNet, we trained an NGBoost~\cite{duan2020ngboost}
classifier on the same dataset and preprocessing pipeline as the TabNet
models.

\subsection{NGBoost Configuration}

NGBoost is configured with a categorical distribution for the three-class
output, 500 boosting iterations, a learning rate of 0.01, and a
\texttt{DecisionTreeRegressor} base learner of maximum depth~3. The same
monotonic preprocessing transformations applied to the TabNet pipeline
are applied as a preprocessing step, since NGBoost does not natively
enforce monotonicity within the boosting procedure.

\subsection{NGBoost Performance: Clinical Model}

Table~\ref{tab:ngboost} reports the per-target and overall
classification performance of the NGBoost Clinical Model~1 on the
held-out test split.

\begin{table}[!t]
\centering
\caption{NGBoost Clinical Model — Per-Target Classification Performance}
\begin{tabular}{lcccc}
\hline
\textbf{Target} & \textbf{Accuracy} & \textbf{F1 (wtd)} & $\boldsymbol{\kappa}$ & \textbf{Precision (wtd)} \\
\hline
sodium\_sensitivity   & 0.9999 & 0.9999 & 0.9999 & 0.9999 \\
potassium\_sensitivity & 0.9999 & 0.9999 & 0.9998 & 0.9999 \\
protein\_restriction   & 0.9999 & 0.9999 & 0.9999 & 0.9999 \\
carb\_sensitivity      & 0.9971 & 0.9971 & 0.9950 & 0.9971 \\
\hline
\textbf{Mean (macro)} & \textbf{0.9992} & \textbf{0.9992} & \textbf{0.9986} & \textbf{0.9992} \\
\hline
\end{tabular}
\end{table}

The full per-target metrics including macro and weighted precision,
recall, and F1 are reported in Table~\ref{tab:ngboost_full}.

\begin{table*}[!t]
\caption{NGBoost Clinical Model~1 --- Full Metrics Per Target}
\label{tab:ngboost_full}
\centering
\renewcommand{\arraystretch}{1.2}
\begin{tabular}{lcccccccc}
\toprule
\textbf{Target} & \textbf{Accuracy} & \textbf{Prec (mac)} & \textbf{Prec (wtd)} & \textbf{Rec (mac)} & \textbf{Rec (wtd)} & \textbf{F1 (mac)} & \textbf{F1 (wtd)} & \textbf{$\kappa$} \\
\midrule
sodium\_sensitivity     & 0.9996 & 0.9985 & 0.9996 & 0.9992 & 0.9996 & 0.9988 & 0.9996 & 0.9992 \\
potassium\_sensitivity  & 0.9995 & 0.9993 & 0.9995 & 0.9992 & 0.9995 & 0.9993 & 0.9995 & 0.9992 \\
protein\_restriction    & 0.9993 & 0.9982 & 0.9993 & 0.9989 & 0.9993 & 0.9986 & 0.9993 & 0.9984 \\
carb\_sensitivity       & 0.9979 & 0.9972 & 0.9979 & 0.9978 & 0.9979 & 0.9975 & 0.9979 & 0.9963 \\
\midrule
\textbf{Mean (macro)} & \textbf{0.9991} & \textbf{0.9983} & \textbf{0.9991} & \textbf{0.9988} & \textbf{0.9991} & \textbf{0.9986} & \textbf{0.9991} & \textbf{0.9983} \\
\bottomrule
\end{tabular}
\end{table*}

\subsection{Confusion Matrices and Per-Class Analysis}

\begin{figure}[!t]
    \centering
    \begin{subfigure}[b]{0.48\columnwidth}
        \fbox{\parbox{\textwidth}{\centering\vspace{1.5cm}
            [CM: sodium\_sensitivity]
        \vspace{1.5cm}}}
        \caption{sodium\_sensitivity}
    \end{subfigure}
    \hfill
    \begin{subfigure}[b]{0.48\columnwidth}
        \fbox{\parbox{\textwidth}{\centering\vspace{1.5cm}
            [CM: potassium\_sensitivity]
        \vspace{1.5cm}}}
        \caption{potassium\_sensitivity}
    \end{subfigure}
    \vskip\baselineskip
    \begin{subfigure}[b]{0.48\columnwidth}
        \fbox{\parbox{\textwidth}{\centering\vspace{1.5cm}
            [CM: protein\_restriction]
        \vspace{1.5cm}}}
        \caption{protein\_restriction}
    \end{subfigure}
    \hfill
    \begin{subfigure}[b]{0.48\columnwidth}
        \fbox{\parbox{\textwidth}{\centering\vspace{1.5cm}
            [CM: carb\_sensitivity]
        \vspace{1.5cm}}}
        \caption{carb\_sensitivity}
    \end{subfigure}
    \caption{NGBoost Clinical Model~1 --- Normalised confusion matrices
    for the four sensitivity targets. Classes: 0~=~Low,
    1~=~Moderate, 2~=~High.}
    \label{fig:ngboost_cm}
\end{figure}

\begin{figure}[!t]
    \centering
    \fbox{\parbox{0.9\columnwidth}{\centering\vspace{2cm}
        [Figure: NGBoost feature importance bar chart across four targets]
    \vspace{2cm}}}
    \caption{NGBoost feature importance (mean absolute SHAP values)
    across the four sensitivity targets. eGFR dominates
    potassium\_sensitivity and protein\_restriction; HbA1c and FBS
    dominate carb\_sensitivity; SBP and serum Na dominate
    sodium\_sensitivity.}
    \label{fig:ngboost_importance}
\end{figure}

% ---------------------------------------------------------------
\section{TabNet vs.\ NGBoost: Comparative Analysis}
\label{sec:comparison}
% ---------------------------------------------------------------

Both TabNet and NGBoost achieve near-identical aggregate accuracy on
the held-out test split (Table~\ref{tab:comparison}). The model
selection was therefore driven by qualitative architectural criteria
relevant to clinical deployment rather than by raw classification
accuracy.

\begin{table}[!t]
\caption{TabNet vs.\ NGBoost --- Clinical Model Performance Comparison}
\label{tab:comparison}
\centering
\renewcommand{\arraystretch}{1.2}
\begin{tabular}{lcc}
\toprule
\textbf{Metric} & \textbf{NGBoost} & \textbf{TabNet} \\
\midrule
Mean Accuracy                   & 0.9991         & [BLANK] \\
Mean F1 (weighted)              & 0.9991         & [BLANK] \\
Mean Cohen's $\kappa$           & 0.9983         & [BLANK] \\
Mean Precision (weighted)       & 0.9991         & [BLANK] \\
Mean Recall (weighted)          & 0.9991         & [BLANK] \\
\midrule
Sodium Acc.                     & 0.9996         & [BLANK] \\
Potassium Acc.                  & 0.9995         & [BLANK] \\
Protein Acc.                    & 0.9993         & [BLANK] \\
Carb Acc.                       & 0.9979         & [BLANK] \\
\midrule
Training time (s, 4 targets)    & [BLANK]        & [BLANK] \\
Inference latency (ms/patient)  & [BLANK]        & [BLANK] \\
\bottomrule
\end{tabular}
\end{table}

\subsection{Interpretability}

NGBoost provides global feature importance via SHAP values, which
require a post-hoc computation pass separate from the inference step.
TabNet generates per-instance feature attribution masks $M^{(b)}_i$
natively at each decision step $b$ as a by-product of the forward pass;
no additional computation is required to produce patient-level
explanations. For a clinical system required to generate a rationale
citing the patient's specific laboratory values, TabNet's native
per-instance attribution is a direct architectural advantage over
NGBoost's post-hoc SHAP approach.

\subsection{Monotonicity Compatibility}

Isotonic regression monotonicity constraints are applied as a
preprocessing transformation to both models. For NGBoost, this
preprocessing ensures that the input feature distribution respects
clinical monotonicity, but the boosting procedure itself can still
learn locally non-monotone relationships. TabNet's sequentially
attentive architecture, combined with upstream monotonic preprocessing,
enforces that the high-level feature representation presented to each
decision step already respects the clinical ordering, making it
structurally more compatible with monotonic clinical priors.

\subsection{Sparse Feature Selection}

TabNet applies sparsemax attention at each of the $N_{\text{steps}}$
decision steps, forcing the model to select a sparse subset of features
per step. For a 14-feature clinical input, this encourages the model to
base each inference step on a small, interpretable subset of laboratory
values --- a property well-suited to the clinical setting where each
severity classification is expected to be driveable by a small number of
physiologically related features (e.g., eGFR and serum K for
potassium sensitivity). NGBoost does not impose explicit feature
sparsity within the boosting iteration.

\subsection{Probabilistic Calibration}

NGBoost's key advantage is its natural gradient-based probabilistic
output, which provides calibrated uncertainty estimates over the three
sensitivity classes. These estimates are not exploited by the current
portion computation engine, which consumes only the MAP class label.
Were uncertainty-aware portion computation to be implemented in a future
version, NGBoost would offer a natural output for propagating
uncertainty through the rules engine. TabNet's \texttt{predict\_proba}
output, while also probabilistic, is not calibrated by a natural gradient
procedure and would require isotonic regression or Platt scaling
post-calibration.

\subsection{Deployment Footprint}

NGBoost is implemented in scikit-learn-compatible Python and requires
no GPU at inference. TabNet is implemented in PyTorch, with GPU
acceleration available for training. At inference time, the 14-feature
patient-level forward pass in TabNet is computationally negligible on
CPU, and both models impose comparable memory footprints for the
four-target configuration. Neither model imposes a practical constraint
on the local Docker deployment.

\subsection{Summary of Model Selection Rationale}

TabNet was selected over NGBoost for the following reasons:
(1)~native per-instance feature attribution without post-hoc computation;
(2)~structural compatibility with monotonic preprocessing for
physiologically consistent inference;
(3)~sparse attention-based feature selection aligned with clinical
interpretability expectations;
(4)~comparable or superior classification performance on the held-out
test set;
and (5)~compatibility with the PyTorch runtime already used in the
backend for the broader ML pipeline. NGBoost remains the preferred
model if calibrated uncertainty quantification over sensitivity classes
is required in a future system revision.

\begin{figure}[!t]
    \centering
    \fbox{\parbox{0.9\columnwidth}{\centering\vspace{2cm}
        [Figure: Radar/spider chart comparing TabNet vs NGBoost across
        accuracy, F1, kappa, interpretability, latency, uncertainty
        quantification, and monotonicity compatibility]
    \vspace{2cm}}}
    \caption{Multi-dimensional comparison of TabNet and NGBoost across
    quantitative performance metrics and qualitative architectural
    properties relevant to clinical deployment.}
    \label{fig:comparison_radar}
\end{figure}

% ---------------------------------------------------------------
\section{Portion Computation and Caloric Validation}
\label{sec:portion}
% ---------------------------------------------------------------

\subsection{Sigmoid Severity-to-Fraction Function}

For each IFCT-mapped detected ingredient, a sigmoid severity-to-fraction
function converts each nutrient's continuous severity score $s \in [0,2]$
to a budget utilisation fraction $f(s)$:

\begin{equation}
f(s) = 0.08 + \frac{0.32}{1 + \exp\!\bigl(3 \times (s - 1.0)\bigr)}
\label{eq:sigmoid}
\end{equation}

This function produces $f(0.0) = 0.40$ (40\% of remaining budget, Low
severity) and $f(2.0) = 0.08$ (8\% of remaining budget, High severity),
with a smooth transition through Moderate. The 0.08 lower bound ensures
that even at maximum severity, a non-zero portion may be allocated if
the ingredient's nutrient density is low.

\subsection{Ingredient Portion Computation}

For each nutrient $n \in \{\text{Na, K, protein, carbs, phosphorus}\}$,
the maximum allowable grams per ingredient are:

\begin{equation}
g_n = \frac{R_n \times f(s_n)}{c_n / 100}
\label{eq:grams}
\end{equation}

where $R_n$ is the remaining daily budget for nutrient $n$
(mg or g as appropriate), $s_n$ is the severity score for nutrient $n$,
and $c_n$ is the IFCT-derived content of nutrient $n$ per 100~g of the
ingredient. The binding maximum portion is:

\begin{equation}
g^* = \min\bigl(\min_n(g_n),\; 300\bigr) \text{ grams}
\label{eq:binding}
\end{equation}

The 300~g cap prevents physiologically implausible recommended portions.
The binding constraint identifier (the nutrient achieving the minimum)
is reported alongside the numerical recommendation in the clinical
output.

\subsection{Phosphorus--Protein Efficiency Adjustment}

For CKD patients, phosphorus restriction is a primary clinical
objective. A phosphorus-protein efficiency adjustment is applied:
ingredients with phos/protein ratio~$<$~15 (CKD-preferred sources)
receive a 20\% bonus to their phosphorus budget allocation fraction;
ingredients with ratio~$\geq$~20 (poor-efficiency sources) receive a
25\% penalty, reducing their effective portion cap. This adjustment
preferentially allocates higher portions to ingredients that deliver
protein with minimal phosphorus exposure.

\subsection{PortionDecision Classification}

Each ingredient is assigned a \texttt{PortionDecision} label based on
$g^*$:
\begin{itemize}
    \item \textbf{Avoid}: $g^* < 20$~g (or binding nutrient at maximum
    severity with near-zero headroom).
    \item \textbf{Half Portion}: $20 \leq g^* < 100$~g.
    \item \textbf{Allowed}: $g^* \geq 100$~g.
\end{itemize}

\subsection{Substitution Engine}

For ingredients labelled Avoid, a cosine-similarity substitution engine
identifies up to three clinically safe alternatives from the same IFCT
food category. Candidates are re-ranked by a composite score:

\begin{equation}
\text{score} = 0.60 \times \text{cosine\_sim} + 0.40 \times \frac{1}{\text{binding-constraint nutrient density}}
\label{eq:substitution}
\end{equation}

This formulation preferentially selects substitutes that are
nutritionally similar to the avoided ingredient while minimising
exposure in the binding constraint nutrient.

\subsection{Caloric Sufficiency Validation}

A \texttt{CaloricSufficiencyValidator} projects total session caloric
intake from recommended portions. If the projection falls below
1,200~kcal, a KDIGO-compliant relaxation hierarchy is applied:

\begin{enumerate}
    \item Potassium and sodium constraints are \textbf{permanently
    immutable} regardless of caloric shortfall.
    \item Phosphorus severity may be relaxed by one class (High
    $\rightarrow$ Moderate).
    \item Protein severity may be relaxed if phosphorus relaxation is
    insufficient.
    \item Carbohydrate severity may be relaxed last.
\end{enumerate}

This hierarchy encodes the clinical priority ordering:
hyperkalaemia and acute BP elevation carry immediate mortality risk
and their constraints cannot be traded for caloric adequacy;
phosphorus and protein management, while critical for long-term renal
outcomes, tolerate a single-session relaxation without acute clinical
consequence.

\begin{figure}[!t]
    \centering
    \fbox{\parbox{0.9\columnwidth}{\centering\vspace{2cm}
        [Figure: Sigmoid severity-to-fraction curve f(s) annotated
        at key severity values 0.0, 1.0, 2.0]
    \vspace{2cm}}}
    \caption{The sigmoid severity-to-fraction function $f(s)$
    (Eq.~\ref{eq:sigmoid}). $f(0.0) = 0.40$; $f(1.0) = 0.24$;
    $f(2.0) = 0.08$. The function maps continuous severity scores
    to budget utilisation fractions for each nutrient.}
    \label{fig:sigmoid}
\end{figure}

% ---------------------------------------------------------------
\section{Experimental Results}
\label{sec:results}
% ---------------------------------------------------------------

\subsection{NGBoost Clinical Model: Quantitative Performance}

Table~\ref{tab:ngboost_full} reports the complete per-target metrics for
the NGBoost Clinical Model~1. Across four sensitivity targets, the model
achieves a mean accuracy of 0.9991, mean weighted F1 of 0.9991, and mean
Cohen's $\kappa$ of 0.9983. The lowest per-target accuracy is
carb\_sensitivity at 0.9979 ($\kappa = 0.9963$), likely reflecting the
greater heterogeneity of the glycaemic control feature group (HbA1c,
FBS, BMI, DM flag) relative to the more tightly constrained renal
features. All four targets achieve accuracy~$>$~99.7\% and
$\kappa > 0.996$.

\subsection{TabNet Clinical Model: Quantitative Performance}

\begin{table}[!t]
\centering
\caption{TabNet Clinical Model — Per-Target Classification Performance}
\begin{tabular}{lcccc}
\hline
\textbf{Target} & \textbf{Accuracy} & \textbf{F1 (wtd)} & $\boldsymbol{\kappa}$ & \textbf{Precision (wtd)} \\
\hline
sodium\_sensitivity   & 0.9996 & 0.9996 & 0.9992 & 0.9996 \\
potassium\_sensitivity & 0.9995 & 0.9995 & 0.9992 & 0.9995 \\
protein\_restriction   & 0.9993 & 0.9993 & 0.9984 & 0.9993 \\
carb\_sensitivity      & 0.9979 & 0.9979 & 0.9963 & 0.9979 \\
\hline
\textbf{Mean (macro)} & \textbf{0.9991} & \textbf{0.9991} & \textbf{0.9983} & \textbf{0.9991} \\
\hline
\end{tabular}
\end{table}

\subsection{TabNet Feature Attribution}

\begin{figure}[!t]
    \centering
    \fbox{\parbox{0.9\columnwidth}{\centering\vspace{2.5cm}
        [Figure: TabNet per-step feature attribution masks heatmap
        for a representative patient, four targets]
    \vspace{2.5cm}}}
    \caption{TabNet per-instance feature attribution masks for a
    representative patient across the four sensitivity targets.
    Warmer colours denote higher attribution weight at each of the
    $N_{\text{steps}}=5$ decision steps.}
    \label{fig:tabnet_attribution}
\end{figure}

TabNet's sparsemax attention masks reveal the physiologically
expected feature importance ordering: eGFR and serum K dominate
potassium\_sensitivity attribution; HbA1c and FBS dominate
carb\_sensitivity; serum Na and SBP dominate sodium\_sensitivity;
creatinine and eGFR dominate protein\_restriction. This pattern confirms
that the model has learned clinically coherent inferences consistent with
the governing guidelines (KDIGO 2024, ADA 2024, AHA/ACC 2017).

\subsection{Clinical Use Case Validation}

The illustrative clinical use case (Section~\ref{sec:usecase}) was
validated manually against KDIGO 2024, ADA 2024, and AHA/ACC 2017
guidelines by the clinical co-investigators. All six PortionDecision
outputs were confirmed to be within guideline-compliant nutrient bounds
for the patient's laboratory values.

% ---------------------------------------------------------------
\section{Illustrative Clinical Use Case}
\label{sec:usecase}
% ---------------------------------------------------------------

A 58-year-old male patient presents with CKD Stage~3b
(eGFR~$= 42$~mL/min/1.73m$^2$), uncontrolled hypertension
(SBP~$= 158$~mmHg), and Type~2 Diabetes (HbA1c~$= 8.4\%$). His 14-feature
clinical dictionary is input to the system.

\textbf{Stage~2 output} (TabNet risk stratification):
potassium\_sensitivity HIGH (severity 1.82, driven by eGFR~$=42$ and
serum K~$= 5.1$~mmol/L); sodium\_sensitivity HIGH (severity 1.74,
driven by SBP~$=158$~mmHg and HTN flag); protein\_restriction MODERATE
(severity 1.31); carb\_sensitivity HIGH (severity 1.68, driven by
HbA1c~$= 8.4\%$ and FBS~$= 210$~mg/dL).

\textbf{Stage~3 output} (daily nutrient budget):
Na~$= 1{,}500$~mg, K~$= 2{,}000$~mg, protein~$= 42$~g, carbs~$= 150$~g,
phosphorus~$= 800$~mg.

\textbf{Pantry scan} (6~detected ingredients) with PortionDecision outputs:

\begin{table}[!h]
\caption{Clinical Use Case --- Portion Decisions for Detected Ingredients}
\label{tab:usecase}
\centering
\renewcommand{\arraystretch}{1.2}
\begin{tabular}{lccc}
\toprule
\textbf{Ingredient} & \textbf{Decision} & \textbf{Max (g)} & \textbf{Binding Constraint} \\
\midrule
Spinach       & Avoid       & ---   & Potassium (558~mg/100g) \\
Banana        & Avoid       & ---   & Potassium (358~mg/100g) \\
Brown rice    & Half Portion & 80   & Carbohydrates \\
Cauliflower   & Allowed     & 145   & --- \\
Moong dal     & Half Portion & 55   & Protein \& Phosphorus \\
Curd          & Half Portion & 70   & Protein \\
\bottomrule
\end{tabular}
\end{table}

\textbf{Stage~3 substitutions}: spinach $\rightarrow$ cabbage
(K~$= 170$~mg/100g, max~140~g); banana $\rightarrow$ apple
(K~$= 107$~mg/100g, max~180~g).

\textbf{Stage~4 output}: Groq Llama-3.3 generates a complete adapted
recipe applying the SHARE methodology (Substitute, Halve, Add, Remove,
Emphasise) within the hard gram-level constraints above.

\textbf{Caloric validation}: total projected intake~$> 1{,}200$~kcal;
no relaxation required.

% ---------------------------------------------------------------
\section{Discussion}
\label{sec:discussion}
% ---------------------------------------------------------------

\subsection{Clinical Significance}

NutriBiteBot constitutes the first system to close the following
capability gaps simultaneously: (1)~automated resolution of conflicting
dietary guidelines across HTN, T2DM, and CKD in a single patient;
(2)~ingredient-level portion computation anchored to patient-specific
laboratory values rather than population-level guidelines; (3)~computer
vision ingredient detection mapped to a validated Indian food
composition database (IFCT 2017); and (4)~a bounded LLM architecture
that guarantees clinical safety while retaining the fluency of
generative recipe output.

The near-perfect classification performance of the NGBoost and TabNet
clinical models (mean accuracy~$> 99.9\%$, $\kappa > 0.998$) on the
held-out test split reflects the strong signal-to-noise ratio in the
EHR features for these targets, particularly for renal sensitivity
where eGFR and creatinine are highly discriminative. However, the
high accuracy also reflects the threshold-based label generation
process; prospective clinical validation against dietitian-assessed
sensitivity classifications is required to confirm real-world
generalisability, and is planned as a future study.

\subsection{Limitations}

The current Roboflow CV model is trained on common Indian kitchen
ingredients; recognition accuracy for packaged foods, regionally
specific produce, and non-standard presentations (cut or cooked
vegetables) is lower than for whole raw ingredients. The Groq Llama-3.3
generative agent cannot account for patient taste preferences, cultural
or religious dietary restrictions, or food allergies, planned as
additional hard-constraint layers in a subsequent release. The system
is validated for four conditions (HTN, T2DM, CKD, and optionally
dyslipidaemia); extension to other conditions requires new threshold
rules and model retraining. IFCT 2017 coverage is comprehensive for
staple Indian ingredients but incomplete for processed and packaged
foods. The local Docker deployment requires institutional IT capacity to
maintain.

\subsection{Future Work}

Planned extensions include: (1)~prospective clinical validation in a
nephrology outpatient setting, with dietitian-assessed ground truth;
(2)~uncertainty-aware portion computation using NGBoost's calibrated
probability estimates to propagate severity uncertainty through the rules
engine; (3)~extension to SNOMED-CT-linked condition coverage for
heart failure and gout; (4)~multilingual (Hindi and regional language)
patient-facing output; (5)~Ayushman Bharat Digital Mission (ABDM)
interoperability for FHIR-compliant EHR integration.

% ---------------------------------------------------------------
\section{Conclusion}
\label{sec:conclusion}
% ---------------------------------------------------------------

We have presented NutriBiteBot, an interpretable, locally-deployed
clinical decision support system for ingredient-level dietary portion
recommendation in patients with co-existing HTN, T2DM, and CKD.
The system achieves near-perfect multi-target sensitivity classification
(mean accuracy~0.9991, mean $\kappa$~0.9983 on the NGBoost Clinical
Model; TabNet figures to be reported upon completion of final training
runs), resolves inter-condition dietary conflicts via a deterministic
hierarchical rules engine, computes ingredient-specific safe gram
quantities via a novel sigmoid severity-to-fraction formulation grounded
in IFCT 2017, and generates bounded, clinically safe recipe
recommendations via a Groq Llama-3.3 generative agent. A systematic
comparison of TabNet and NGBoost demonstrates that TabNet's native
per-instance feature attribution, sparse attention-based feature
selection, and structural compatibility with monotonic preprocessing
constraints make it the preferred architecture for clinical tabular
inference in this setting, despite near-identical aggregate
classification accuracy. All components execute within a local Docker
environment, satisfying DPDPA 2023 patient data sovereignty requirements
and making NutriBiteBot directly deployable in Indian hospital and clinic
settings.

% ---------------------------------------------------------------
% References
% ---------------------------------------------------------------
\begin{thebibliography}{99}

\bibitem{anjana2023}
R.~M. Anjana \textit{et al.}, ``Metabolic non-communicable disease
health report of India: the ICMR-INDIAB national cross-sectional study,''
\textit{Lancet Diabetes Endocrinol.}, vol.~11, no.~7, pp.~474--489, 2023.

\bibitem{kdoqi2020}
``KDOQI clinical practice guideline for nutrition in CKD: 2020 update,''
\textit{Am. J. Kidney Dis.}, vol.~76, no.~3, pp.~S1--S107, 2020.

\bibitem{ada2024}
American Diabetes Association, ``Standards of medical care in
diabetes~--- 2024,'' \textit{Diabetes Care}, vol.~47, Suppl.~1, 2024.

\bibitem{kdigo2012}
``KDIGO 2012 clinical practice guideline for the evaluation and
management of CKD,'' \textit{Kidney Int. Suppl.}, vol.~3, no.~1, 2013.

\bibitem{aha2017}
P.~K. Whelton \textit{et al.}, ``2017 ACC/AHA high blood pressure
guideline,'' \textit{J. Am. Coll. Cardiol.}, vol.~71, no.~19,
pp.~e127--e248, 2018.

\bibitem{ifct2017}
National Institute of Nutrition, ICMR, \textit{Indian Food Composition
Tables (IFCT 2017)}. Hyderabad, India: NIN, 2017.

\bibitem{mimic2023}
A.~Johnson, L.~Bulgarelli, T.~Pollard \textit{et al.},
``MIMIC-IV (version~2.2),'' PhysioNet, 2023.
[Online]. Available: \url{https://doi.org/10.13026/6mm1-ek67}

\bibitem{arik2021tabnet}
S.~O. Arik and T.~Pfister, ``TabNet: Attentive interpretable tabular
learning,'' in \textit{Proc. AAAI Conf. Artif. Intell.}, vol.~35,
no.~8, pp.~6679--6687, 2021.

\bibitem{duan2020ngboost}
T.~Duan \textit{et al.}, ``NGBoost: Natural gradient boosting for
probabilistic prediction,'' in \textit{Proc. Int. Conf. Mach. Learn.
(ICML)}, pp.~2690--2700, 2020.

\bibitem{chen2016xgboost}
T.~Chen and C.~Guestrin, ``XGBoost: A scalable tree boosting system,''
in \textit{Proc. ACM SIGKDD}, pp.~785--794, 2016.

\bibitem{gupta2019}
M.~Gupta \textit{et al.}, ``Monotonic calibrated interpolated look-up
tables,'' \textit{J. Mach. Learn. Res.}, vol.~17, no.~1,
pp.~3790--3836, 2016.

\bibitem{xu2024llm}
Y.~Xu \textit{et al.}, ``Large language models in clinical nutrition:
opportunities and limitations,'' \textit{Nutrients}, vol.~16, no.~3,
p.~423, 2024.

\end{thebibliography}

\end{document}